\documentclass[11pt,a4paper]{article}
\usepackage[T1]{fontenc}
\usepackage[utf8]{inputenc}
\usepackage{lmodern,amsmath,amssymb}
\usepackage[margin=25mm]{geometry}
\usepackage{longtable,booktabs,array,calc}
\usepackage[hidelinks]{hyperref}
\hypersetup{pdftitle={The large-field action lower bound is immediate in the natural variable, and the constant },pdfauthor={navstokgap research working notes}}
\newcommand{\tightlist}{\setlength{\itemsep}{0pt}\setlength{\parskip}{0pt}}
\setlength{\emergencystretch}{3em}
\setcounter{secnumdepth}{0}
\begin{document}
\hypertarget{the-large-field-action-lower-bound-is-immediate-in-the-natural-variable-and-the-constant-field-saturates-it}{%
\section{The large-field action lower bound is immediate in the natural
variable, and the constant field saturates
it}\label{the-large-field-action-lower-bound-is-immediate-in-the-natural-variable-and-the-constant-field-saturates-it}}

The bound asked for at the end of
\href{large-field-entropy-count.md}{the entropy note} holds with the
optimal constant, once the large-field condition is written in the
variable the action already measures. Define the region as
\[K_\eta=\Big\{A:\ \frac1{\ell^4}\int_{\rm block}\big|G_t(x)\big|^2d^4x\ \ge\ \Big(\frac{\eta}{\ell^2}\Big)^2\Big\},
\qquad \sqrt{8t}=\ell,\] with \(G_t\) the field strength at flow time
\(t\). Then
\[\inf_{K_\eta}\frac{S_E}{\hbar}\ \ge\ \frac{\eta^2}{4g^2},\] by two
steps only: the gradient flow decreases the Euclidean action, so
\(S_E[A]\ge S_E[\text{flowed }A]\), and the flowed action is at least
its restriction to the block, which the constraint bounds below. The
bound is \textbf{uniform in the block size, in the lattice spacing, in
the volume and in the gauge group} in the normalization used here, and
it is \textbf{saturated}: a field of magnitude \(\eta/\ell^2\) supported
on the block attains it, so the constant \(1/4\) is optimal. The
Nielsen--Olesen instability is no obstruction, because every deformation
that lowers the action also lowers the flowed block average and
therefore leaves \(K_\eta\): the constant field is a saddle of the
unconstrained action and a minimizer of the constrained one. Together
with \href{large-field-entropy-count.md}{the entropy count} this closes
the large-field corner as an estimate problem: cost \(\eta^2/(4g^2)\),
count \(\eta^2/(8\pi^2)\), net exponent
\(-\frac{\eta^2}{4g^2}\big(1-Cg^2/2\pi^2\big)\). Constants explicit;
nothing promoted.

\hypertarget{the-constraint-in-the-right-variable}{%
\subsection{1. The constraint in the right
variable}\label{the-constraint-in-the-right-variable}}

The Euclidean action is the \(L^2\) norm of the field strength,
\[\frac{S_E}{\hbar}=\frac1{4g^2}\int\big|F^a_{\mu\nu}(x)\big|^2d^4x ,\]
so the natural measure of ``how large the field is in a block'' is the
same quantity restricted to the block, and the natural smoothing is the
gradient flow at the radius of the block. Take \(\sqrt{8t}=\ell\) and
\[\mathcal F_\ell(A)=\frac1{\ell^4}\int_{\rm block}\big|G_t(x)\big|^2d^4x,
\qquad K_\eta=\Big\{A:\ \mathcal F_\ell(A)\ge\eta^2/\ell^4\Big\}.\] This
is gauge invariant, because \(|G_t|^2\) is, and it is the quantity that
the truncation estimate of \href{flow-jacobian-truncation-error.md}{the
Jacobian note} is sensitive to, up to replacing a supremum by a block
average.

\hypertarget{the-bound}{%
\subsection{2. The bound}\label{the-bound}}

\textbf{Proposition 1.} For every \(\ell>0\), every lattice spacing
\(a<\ell\) and every volume,
\[\inf_{A\in K_\eta}\ \frac{S_E[A]}{\hbar}\ \ge\ \frac{\eta^2}{4g^2}.\]

\emph{Proof.} Two steps.

\emph{(i) The flow decreases the action.} Along the gradient flow
\(\partial_sB=-D^*G\) one has \(\frac{d}{ds}S_E(B_s)=-\|D^*G\|_2^2\le0\)
(\href{comparison-and-bridges.md}{comparison note} §4; Lüscher,
arXiv:1006.4518v3, after equation (1.4), passage level), so
\(S_E[A]=S_E[B_0]\ge S_E[B_t]=\frac{\hbar}{4g^2}\int|G_t|^2d^4x\).

\emph{(ii) Positivity outside the block.} The integrand is nonnegative,
so
\[\frac1{4g^2}\int\big|G_t\big|^2\ \ge\ \frac1{4g^2}\int_{\rm block}\big|G_t\big|^2
\ \ge\ \frac1{4g^2}\,\ell^4\cdot\frac{\eta^2}{\ell^4}=\frac{\eta^2}{4g^2},\]
the middle inequality being the definition of \(K_\eta\). \(\square\)

No property of the gauge group beyond the normalization of the trace
enters, and no property of the lattice beyond \(a<\ell\), which is
needed only so that the flow at radius \(\ell\) is meaningful.

\textbf{Proposition 2 (sharpness).} The constant \(1/4\) cannot be
improved. A configuration whose field strength equals \(\eta/\ell^2\) on
the block and vanishes outside has \(S_E/\hbar=\eta^2/(4g^2)\) and lies
in \(K_\eta\) up to the smoothing at the boundary of the block, which
contributes a relative correction of order \(a/\ell\).

\hypertarget{why-the-instability-does-not-obstruct-the-bound}{%
\subsection{3. Why the instability does not obstruct the
bound}\label{why-the-instability-does-not-obstruct-the-bound}}

\href{flow-instability-large-field.md}{The instability note} shows that
a constant chromomagnetic field is a saddle of the unconstrained action,
with \(\mathcal N(\eta)=\eta^2/(8\pi^2)\) descent directions. Each of
those directions lowers \(S_E\), and by step (i) of Proposition 1 it
therefore lowers \(\int|G_t|^2\) as well, so it lowers
\(\mathcal F_\ell\) and moves the configuration out of \(K_\eta\). The
two statements are consistent and say different things:

\begin{itemize}
\tightlist
\item
  \emph{unconstrained}: the constant field is not a local minimum, and
  the flow runs away from it at rate \(2\|G\|\);
\item
  \emph{constrained}: on \(K_\eta\) the constant field attains the
  infimum, because every descent direction violates the constraint.
\end{itemize}

The instability therefore affects the \textbf{treatment} of the region,
which must expand around an inhomogeneous configuration if one insists
on staying inside \(K_\eta\) while following the flow, and not the
\textbf{cost} of the region, which Proposition 1 fixes.

\hypertarget{the-corner-assembled}{%
\subsection{4. The corner, assembled}\label{the-corner-assembled}}

Collecting the three statements about a large-field region of strength
\(\eta\) at scale \(\ell\):

\begin{longtable}[]{@{}
  >{\raggedright\arraybackslash}p{(\columnwidth - 4\tabcolsep) * \real{0.3333}}
  >{\raggedright\arraybackslash}p{(\columnwidth - 4\tabcolsep) * \real{0.3333}}
  >{\raggedright\arraybackslash}p{(\columnwidth - 4\tabcolsep) * \real{0.3333}}@{}}
\toprule\noalign{}
\begin{minipage}[b]{\linewidth}\raggedright
quantity
\end{minipage} & \begin{minipage}[b]{\linewidth}\raggedright
value
\end{minipage} & \begin{minipage}[b]{\linewidth}\raggedright
source
\end{minipage} \\
\midrule\noalign{}
\endhead
\bottomrule\noalign{}
\endlastfoot
minimal action cost & \(\eta^2/(4g^2)\), sharp & Propositions 1--2
here \\
unstable directions & \(\eta^2/(8\pi^2)\) &
\href{large-field-entropy-count.md}{entropy note} §2 \\
flow-truncation error inside & \(e^{2\eta}\), sharp &
\href{flow-instability-large-field.md}{instability note} \\
net weight &
\(\exp\big[-\frac{\eta^2}{4g^2}\big(1-\frac{Cg^2}{2\pi^2}\big)\big]\) &
combining the first two \\
\end{longtable}

All four are independent of \(\ell\) and of the lattice spacing, which
is the scale invariance that makes an induction over scales possible at
all. The first three are proved or sharp; the fourth needs the standard
conditional argument comparing the constrained partition function with
the full one, in which the entropy of the block is exactly the count in
the second row.

\textbf{What is closed.} The large-field region is no longer an
unquantified obstacle. Its cost, its entropy and the failure mode of the
flow inside it are all explicit, scale-invariant, and consistent with
each other.

\textbf{What is not.} How the effective Hamiltonian is defined inside
the region. Proposition 1 bounds the weight of the region; it does not
say what replaces the flow-and-truncate step there, and the constructive
answer, an expansion around the constrained minimizer with its
\(\eta^2/(8\pi^2)\) soft directions treated separately, is a
construction rather than an estimate.

\hypertarget{consequence-for-state}{%
\subsection{5. Consequence for STATE}\label{consequence-for-state}}

The lower bound asked for is proved with the optimal constant and
uniformly in every parameter, by flow monotonicity and positivity, once
the constraint is written as a block average of the flowed action
density. The large-field corner is closed at the level of estimates. The
remaining content of the decimation step is the construction inside the
region, which is the same constructive problem the programme has reached
from three directions now: after the small-field truncation is
controlled, the large-field regions are rare and expensive but must
still be given an effective description.
\end{document}
